%% toc: P15 163 163.1,163.2,163.3,163.4
% Proofs for the boundaries of the exchangeable fields, averages, directing laws and finite-urn bound.
\section{Boundaries of the exchangeable fields, averages, directing laws and finite-urn bound: proofs}
\label{sec:annex-p15-w02c-p15-022}

Throughout, \(\Omega=S^{\mathbb N_+}\) carries the product sigma-field
\(\mathcal F\), the coordinates \(X_j(\omega)=\omega_j\), and the finite
permutations \((T_\pi\omega)_j=\omega_{\pi(j)}\) of
\cref{def:p15-permutable-events}. Two facts are used repeatedly.

\emph{(F1) Events that ignore the first coordinate.} Let
\(\mathcal L\) be the class of events \(C\in\mathcal F\) such that
\(\omega\in C\Leftrightarrow\omega'\in C\) whenever
\(\omega_j=\omega'_j\) for every \(j>1\). It is closed under complements
and countable unions, so it is a sigma-field, and it contains every
\(X_j^{-1}(G)\) with \(j>1\) and \(G\in\mathcal S\). Hence
\(\mathcal T\subseteq\sigma(X_j:j>1)\subseteq\mathcal L\).

\emph{(F2) Transport by a permutation.} If \(\pi\) is a permutation of
\(\{1,\ldots,n\}\) extended by the identity, then
\(X_j\circ T_\pi=X_{\pi(j)}\) and
\(\mathbf1_C\circ T_\pi=\mathbf1_{T_\pi^{-1}C}\). If \(P\) is
exchangeable, \(P\circ T_\pi^{-1}=P\) by
\cref{thm:p15-exchangeable-canonical}, so
\(E[F\circ T_\pi]=E[F]\) for every measurable \(F:\Omega\to[0,\infty]\),
and for every \(F\) for which either side is finite.

\begin{annexitem}{annex:p15-w02c-p15-022-field-strictness}{thm:p15-w02c-p15-022-field-strictness}{When the invariant inclusions are strict}
\begin{proof}
\emph{Item 1.} Each cylinder \(\{X_1\in B_1,\ldots,X_m\in B_m\}\) with
\(B_j\in\{\varnothing,S\}\) is \(\varnothing\) or \(\Omega\). The class
\(\{\varnothing,\Omega\}\) is a sigma-field containing these
generators, so \(\mathcal F=\{\varnothing,\Omega\}\). The fields
\(\mathcal E_n\), \(\mathcal E\) and \(\mathcal T\) are sub-sigma-fields
of \(\mathcal F\), so each equals \(\{\varnothing,\Omega\}\).

\emph{Item 2.} Choose \(x\in B\) and \(y\in S\setminus B\). The event
\(A_n=\bigcap_{j=1}^nX_j^{-1}(B)\) is measurable. If \(\pi\) permutes
\(\{1,\ldots,n\}\) and fixes every larger index, then
\(T_\pi\omega\in A_n\) means \(\omega_{\pi(j)}\in B\) for \(1\leq j\leq n\),
which is the condition \(\omega\in A_n\) because \(\pi\) maps
\(\{1,\ldots,n\}\) onto itself. So \(T_\pi^{-1}A_n=A_n\) and
\(A_n\in\mathcal E_n\). Let \(\tau\) be the transposition of \(n\) and
\(n+1\); it permutes \(\{1,\ldots,n+1\}\). The point
\(\omega=(x,\ldots,x,y,y,\ldots)\) with \(n\) leading entries \(x\) lies
in \(A_n\), and \((T_\tau\omega)_n=\omega_{n+1}=y\notin B\), so
\(\omega\notin T_\tau^{-1}A_n\). Thus \(T_\tau^{-1}A_n\neq A_n\) and
\(A_n\notin\mathcal E_{n+1}\).

The event \(A\) is the countable union over \(i\) of
\(X_i^{-1}(B)\cap\bigcap_{j\neq i}X_j^{-1}(S\setminus B)\), so it is
measurable. For a finite permutation \(\pi\), the map \(j\mapsto\pi(j)\)
is a bijection of \(\mathbb N_+\) carrying
\(\{j:\omega_{\pi(j)}\in B\}\) onto \(\{i:\omega_i\in B\}\); the two sets
have the same number of elements. So \(T_\pi\omega\in A\) exactly when
\(\omega\in A\), and \(A\in\mathcal E\). The points
\(a=(x,y,y,\ldots)\) and \(b=(y,y,y,\ldots)\) agree beyond the first
coordinate, \(a\in A\) and \(b\notin A\). Hence \(A\notin\mathcal L\), and
by (F1) \(A\notin\mathcal T\).

\emph{Item 3.} Suppose \(A_n\in\mathcal E_{n+1}^{P}\): there is
\(C\in\mathcal E_{n+1}\) with \(P(A_n\mathbin{\triangle}C)=0\). With
\(\tau\) as above, \(T_\tau^{-1}C=C\) and
\(T_\tau^{-1}(A_n\mathbin{\triangle}C)
 =T_\tau^{-1}A_n\mathbin{\triangle}C\). By (F2),
\(P(T_\tau^{-1}A_n\mathbin{\triangle}C)=P(A_n\mathbin{\triangle}C)=0\).
Since \(A_n\mathbin{\triangle}T_\tau^{-1}A_n\) is contained in the union
of \(A_n\mathbin{\triangle}C\) and \(C\mathbin{\triangle}T_\tau^{-1}A_n\),
it is null. But it contains
\(G=\{X_1\in B,\ldots,X_n\in B,\ X_{n+1}\notin B\}\): a point of \(G\)
is in \(A_n\), and its image under \(T_\tau\) has \(n\)th coordinate
outside \(B\). Since \(P(G)>0\), this is a contradiction.

Under an i.i.d.\ law with coordinate law \(\mu\) and
\(0<\mu(B)<1\), \(P(G)=\mu(B)^n(1-\mu(B))>0\).

If \(P\) is the law of \((\xi,\xi,\ldots)\), then \(A_{n+1}\in\mathcal
E_{n+1}\) by the first part of item~2 at level \(n+1\), and
\(A_n\mathbin{\triangle}A_{n+1}=A_n\setminus A_{n+1}=G\). Its probability
is the probability that \(\xi\in B\) and \(\xi\notin B\), namely zero.
Hence \(A_n\in\mathcal E_{n+1}^{P}\).

\emph{Item 4.} Under this \(P\), \(P(X_1=1)=1/2\) and \(P(X_2=1)=0\), so
the vector \((X_1,X_2)\) and its permutation \((X_2,X_1)\) have different
laws, and \(P\) is not exchangeable. The points \(a,b\) agree beyond the
first coordinate. By (F1), every \(C\in\mathcal T\) contains both or
neither. If both, \(P(C)=1\) and \(A\mathbin{\triangle}C\supseteq\{b\}\)
because \(b\notin A\); if neither, \(P(C)=0\) and
\(A\mathbin{\triangle}C\supseteq\{a\}\) because \(a\in A\). In either case
\(P(A\mathbin{\triangle}C)\geq1/2\). So \(A\notin\mathcal T^{P}\), while
\(A\in\mathcal E\subseteq\mathcal E^{P}\) by item~2, and
\(P(A)=P(\{a\})=1/2\).
\end{proof}
\end{annexitem}

\begin{annexitem}{annex:p15-w02c-p15-022-partial}{thm:p15-w02c-p15-022-not-exchangeable}{Partial symmetry and the canonical reduction}
\begin{proof}
\emph{Item 1.} The law \(P\) gives probability \(1/2\) to each of
\(c=(1,0,1,0,\ldots)\) and \(c'=(0,1,0,1,\ldots)\). Each coordinate is
\(1\) at exactly one of them, so every \(X_j\) is a fair bit. The pair
\((X_1,X_2)\) takes the values \((1,0)\) and \((0,1)\) with probability
\(1/2\) each, and so does \((X_2,X_1)\). The triple \((X_1,X_2,X_3)\)
takes the values \((1,0,1)\) and \((0,1,0)\), while \((X_1,X_3,X_2)\)
takes \((1,1,0)\) and \((0,0,1)\); these laws are different, so \(P\) is
not exchangeable. By (F2), \(T_\tau^{-1}\{X_1=X_3\}=\{X_2=X_3\}\). At
both \(c\) and \(c'\), \(X_1=X_3\) and \(X_2\neq X_3\), so
\(P(X_1=X_3)=1\) and \(P(X_2=X_3)=0\).

\emph{Item 2.} Both \(Q_1\) and \(Q_2\) give \(Z\) the law \(P\): for
\(Q_1\) because \(Z\) is the first factor of a product law, for \(Q_2\)
because \(Z\) composed with \(\omega\mapsto(\omega,\omega_1)\) is the
identity. Every event of \(Z^{-1}(\mathcal E)\) has the form
\(Z^{-1}(A_0)\) with \(A_0\in\mathcal E\).

Under \(Q_1\), by the product structure,
\[
 E_{Q_1}\bigl[U\mathbf1_{Z^{-1}(A_0)}\bigr]
   =\tfrac12P(A_0)
   =E_{Q_1}\bigl[\tfrac12\mathbf1_{Z^{-1}(A_0)}\bigr],
\]
and the constant \(1/2\) is measurable with respect to every
sigma-field. So it is a version of \(E_{Q_1}(U\mid Z^{-1}(\mathcal E))\).

Under \(Q_2\), the event \(\{U=X_1\circ Z\}\) is measurable and
contains the image of \(\omega\mapsto(\omega,\omega_1)\), so it has
probability one, and
\(E_{Q_2}[U\mathbf1_{Z^{-1}(A_0)}]=E_P[X_1\mathbf1_{A_0}]\). The law
\(P\) gives probability \(1/2\) to each of \((1,1,\ldots)\) and
\((0,0,\ldots)\), and at both points \(X_1=\bar Y\); the measurable event
\(\{X_1\neq\bar Y\}\) contains neither and is null. For each \(m\),
\(\bar Y=\limsup_n(S_n-S_m)/n\) because \(0\leq S_m/n\leq m/n\to0\), and
the right side is \(\sigma(X_j:j>m)\)-measurable. So \(\bar Y\) is
\(\mathcal T\)-measurable, hence \(\mathcal E\)-measurable by
\cref{thm:p15-exchangeable-canonical}. Therefore
\[
 E_{Q_2}\bigl[U\mathbf1_{Z^{-1}(A_0)}\bigr]
   =E_P[\bar Y\mathbf1_{A_0}]
   =E_{Q_2}\bigl[(\bar Y\circ Z)\mathbf1_{Z^{-1}(A_0)}\bigr],
\]
and \(\bar Y\circ Z\) is \(Z^{-1}(\mathcal E)\)-measurable. It is a
version of \(E_{Q_2}(U\mid Z^{-1}(\mathcal E))\), and it equals \(1\) or
\(0\) with \(Q_2\)-probability \(P(\bar Y=1)=P(\bar Y=0)=1/2\).
\end{proof}
\end{annexitem}

\begin{annexitem}{annex:p15-w02c-p15-022-average-edges}{thm:p15-w02c-p15-022-average-edges}{The index range and the distinctness of indices}
\begin{proof}
\emph{Item 1.} If \(n<k\), no \(k\) distinct indices lie in a set of
\(n\) elements, so \(I_{n,k}=\varnothing\); the product
\(n(n-1)\cdots(n-k+1)\) contains the factor \(n-n=0\). At \(n=k\), a
tuple of \(k\) distinct indices from \(\{1,\ldots,k\}\) is an ordering
\((\pi(1),\ldots,\pi(k))\) of them, and \((k)_k=k!\).

\emph{Item 2.} Let \(n\geq2\). The map \((i_1,i_2)\mapsto(i_2,i_1)\) is a
bijection of \(I_{n,2}\) onto itself without fixed points, because the
indices are distinct. It partitions \(I_{n,2}\) into pairs whose two
summands are \(\varphi(X_{i_1},X_{i_2})\) and
\(\varphi(X_{i_2},X_{i_1})=-\varphi(X_{i_1},X_{i_2})\). So the sum, and
\(A_{n,2}(\varphi)\), vanish at every point. By
\cref{thm:p15-exchangeable-averages},
\(E(\varphi(X_1,X_2)\mid\mathcal E_n)=A_{n,2}(\varphi)=0\) almost surely.
For \(\varphi(x,y)=x-y\), \(X_1\) and \(X_2\) have the same law, so
\(X_1-X_2\) is integrable, and linearity of conditional expectation
gives \(E(X_1\mid\mathcal E_n)=E(X_2\mid\mathcal E_n)\) almost surely.

\emph{Item 3.} Split the pairs \((i,j)\in\{1,\ldots,n\}^2\) into those
with \(i\neq j\), which form \(I_{n,2}\), and the diagonal \(i=j\):
\[
 \sum_{i,j=1}^n\varphi(X_i,X_j)
   =(n)_2A_{n,2}(\varphi)+nA_{n,1}(\varphi_\Delta).
\]
Divide by \(n^2\) and use \((n)_2=n(n-1)\). If \(\varphi(X_1,X_2)\) and
\(\varphi_\Delta(X_1)=\varphi(X_1,X_1)\) are integrable,
\cref{thm:p15-exchangeable-averages} with \(k=2\) and with \(k=1\) gives
\(A_{n,2}(\varphi)=E(\varphi(X_1,X_2)\mid\mathcal E_n)\) and
\(A_{n,1}(\varphi_\Delta)=E(\varphi(X_1,X_1)\mid\mathcal E_n)\), and
linearity gives the displayed conditional expectation. Finally
\[
 V_n(\varphi)-A_{n,2}(\varphi)
   =\frac1n\bigl(A_{n,1}(\varphi_\Delta)-A_{n,2}(\varphi)\bigr).
\]
By the same theorem both averages converge almost surely to finite
limits, so the bracket is almost surely bounded in \(n\), and the
difference tends to zero almost surely.

\emph{Item 4.} Here \(I_{2,2}=\{(1,2),(2,1)\}\), so
\(A_{2,2}(\varphi)=\tfrac12(X_1X_2+X_2X_1)=X_1X_2\), which equals
\(E(X_1X_2\mid\mathcal E_2)\) by
\cref{thm:p15-exchangeable-averages}. Also
\(V_2(\varphi)=\tfrac14\sum_{i,j\leq2}X_iX_j=(X_1+X_2)^2/4\). On
\(\{X_1+X_2=1\}\), \(X_1X_2=0\) and \(V_2(\varphi)=1/4\); this event has
probability \(P(X_1=1,X_2=0)+P(X_1=0,X_2=1)=1/2\). For general \(n\),
\(X_i^2=X_i\) gives
\(\sum_{i\neq j}X_iX_j=S_n^2-\sum_iX_i^2=S_n(S_n-1)\), so
\(A_{n,2}(\varphi)=S_n(S_n-1)/(n(n-1))\), and
\(\sum_{i,j}X_iX_j=S_n^2\) gives \(V_n(\varphi)=S_n^2/n^2\).
\end{proof}
\end{annexitem}

\begin{annexitem}{annex:p15-w02c-p15-022-integrability}{thm:p15-w02c-p15-022-integrability}{The integrability hypothesis and the exceptional null sets}
\begin{proof}
\emph{Item 1.} For any measurable \(\psi:S^k\to\mathbb R\), a permutation
of \(\{1,\ldots,n\}\) permutes the summands of \(A_{n,k}(\psi)\), as in
\cref{annex:p15.exchangeable-averages}; so \(A_{n,k}(\psi)\) is
\(\mathcal E_n\)-measurable. In particular \(W_n\) is
\(\mathcal E_n\)-measurable and finite at every point, being a finite
average of finite numbers. Hence \(D_c\in\mathcal E_n\), and \(D_c\)
increases to \(\Omega\) as \(c\to\infty\).

Let \(\psi\geq0\) be measurable, \(D\in\mathcal E_n\), and
\(i\in I_{n,k}\). Choose a permutation \(\pi\) of \(\{1,\ldots,n\}\) with
\(\pi(j)=i_j\) for \(1\leq j\leq k\). By (F2),
\(\psi(X_{i_1},\ldots,X_{i_k})\mathbf1_D
 =(\psi(X_1,\ldots,X_k)\mathbf1_D)\circ T_\pi\), since
\(T_\pi^{-1}D=D\), and
\[
 E[\psi(X_{i_1},\ldots,X_{i_k})\mathbf1_D]
   =E[\psi(X_1,\ldots,X_k)\mathbf1_D]\in[0,\infty].
\tag{\(*\)}
\]
Averaging over \(I_{n,k}\) gives
\(E[A_{n,k}(\psi)\mathbf1_D]=E[\psi(X_1,\ldots,X_k)\mathbf1_D]\). With
\(\psi=|\varphi|\) and \(D=D_c\),
\(E[|\varphi(X_1,\ldots,X_k)|;D_c]=E[W_n;D_c]\leq c\). By \((*)\) with
\(\psi=|\varphi|\), every summand \(\varphi(X_{i_1},\ldots,X_{i_k})\)
is integrable on \(D_c\), with the same bound.

Now let \(D\in\mathcal E_n\). Then \(D\cap D_c\in\mathcal E_n\), and
\((*)\) applied to \(\varphi^+\) and to \(\varphi^-\) on \(D\cap D_c\)
gives two finite equalities. Subtracting,
\(E[\varphi(X_{i_1},\ldots,X_{i_k})\mathbf1_{D\cap D_c}]
 =E[\varphi(X_1,\ldots,X_k)\mathbf1_{D\cap D_c}]\), and averaging over
\(I_{n,k}\),
\[
 E\bigl[\mathbf1_{D_c}A_{n,k}(\varphi)\,\mathbf1_D\bigr]
   =E\bigl[\mathbf1_{D_c}\varphi(X_1,\ldots,X_k)\,\mathbf1_D\bigr].
\]
The variable \(\mathbf1_{D_c}A_{n,k}(\varphi)\) is integrable, being
bounded in absolute value by \(\mathbf1_{D_c}W_n\leq c\), and
\(\mathcal E_n\)-measurable. Since the identity holds for every
\(D\in\mathcal E_n\), it is a version of
\(E(\mathbf1_{D_c}\varphi(X_1,\ldots,X_k)\mid\mathcal E_n)\).

\emph{Item 2.} The probabilities sum to
\(\sum_{m\geq1}2\cdot2^{-m-1}=1\), and
\(E|\xi|=\sum_{m\geq1}2^m\cdot2\cdot2^{-m-1}=\infty\). For cylinders,
\(P(X_1\in B_1,\ldots,X_r\in B_r)=P(\xi\in B_1\cap\cdots\cap B_r)\),
which does not change when the \(B_j\) are permuted; so every finite
coordinate vector is exchangeable and \(P\) is exchangeable. The
variable \(\varphi(X_1,\ldots,X_k)=X_1\) has the law of \(\xi\), so it is
not integrable. Each \(m\in\{1,\ldots,n\}\) is the first entry of
exactly \((n-1)(n-2)\cdots(n-k+1)\) tuples of \(I_{n,k}\), the number of
ways to choose the remaining \(k-1\) distinct entries in order; this
count is \(1\) when \(k=1\). Since
\((n)_k=n\cdot(n-1)\cdots(n-k+1)\), it follows that
\(A_{n,k}(\varphi)=n^{-1}\sum_{m=1}^nX_m=S_n/n\) at every point. The
event \(\bigcap_{j}\{X_j=X_1\}\) is measurable and has probability one,
because each \(\{X_j\neq X_1\}\) has probability
\(P(\xi\neq\xi)=0\). On it \(S_n/n=X_1\) for every \(n\), so the
averages converge almost surely to \(X_1\). Since
\(E|A_{n,k}(\varphi)|=E|X_1|=E|\xi|=\infty\), neither the averages nor
the limit is integrable; \(E(X_1\mid\mathcal E_n)\) and
\(E(X_1\mid\mathcal E)\) are not defined for this variable, since
\(E\xi^+=E\xi^-=\infty\), and there is no \(L^1\) convergence to assert.

\emph{Item 3.} The formula \(A_{n,k}(\varphi)=S_n/n\) was proved in
item~2 for every law. Let \(Y=|X_1|\). Since
\(EY=\int_0^\infty P(Y>t)\,dt\), by Fubini's theorem applied to
\((\omega,t)\mapsto\mathbf1_{\{Y(\omega)>t\}}\), and \(t\mapsto P(Y>t)\) is decreasing,
\[
 \sum_{n\geq1}P(Y>n)\geq\sum_{n\geq1}\int_n^{n+1}P(Y>t)\,dt
   =\int_1^\infty P(Y>t)\,dt\geq EY-1=\infty.
\]
The events \(\{|X_n|>n\}\) are mutually independent and
\(P(|X_n|>n)=P(Y>n)\), so \cref{thm:p15-borel-cantelli-second} gives
\(P(|X_n|>n\text{ i.o.})=1\). On the event that \(S_n/n\to L\) for a
finite \(L\),
\[
 \frac{X_n}n=\frac{S_n}n-\frac{n-1}n\cdot\frac{S_{n-1}}{n-1}\longrightarrow L-L=0,
\]
so \(|X_n|>n\) for only finitely many \(n\). That event is therefore
disjoint from an event of probability one, and has probability zero.
For the law of \(\xi\) in item~2, \(E|X_1|=E|\xi|=\infty\).

\emph{Item 4.} Countable subsets of \([0,1]\) are Borel. Let \(\lambda\) be
Lebesgue measure. For countable \(D\), countable additivity gives
\[
 P(X_1\in D)=\sum_{x\in D}\lambda(\{x\})=0,
\]
so \(P(X_j\in D\text{ for some }j)=0\). Off this null
event \(A_{n,1}(\mathbf1_D)=n^{-1}\#\{j\leq n:X_j\in D\}=0\) for every
\(n\). Since \(\mathbf1_D(X_1)=0\) almost surely, the constant \(0\) is a
version of \(E(\mathbf1_D(X_1)\mid\mathcal E)\).

For \(i\neq j\), \((X_i,X_j)\) has law \(\lambda\otimes\lambda\), and the
diagonal \(\{(s,t):s=t\}\) is closed. By Fubini's theorem its measure is
\(\int_0^1\lambda(\{s\})\,ds=0\). So
\(N_0=\bigcup_{i\neq j}\{X_i=X_j\}\) is a null event. Let
\[
 J=\bigcup_{m\geq0}\{j\in\mathbb N_+:4^m\leq j<2\cdot4^m\},
 \qquad c(n)=\#(J\cap\{1,\ldots,n\}).
\]
The pieces of \(J\) have \(4^m\) elements, and no element of \(J\) lies in
\([2\cdot4^m,4^{m+1})\). Hence
\(c(2\cdot4^m-1)=c(4^{m+1}-1)=\sum_{i=0}^m4^i=(4^{m+1}-1)/3\), so
\[
 \frac{c(4^{m+1}-1)}{4^{m+1}-1}=\frac13,
 \qquad
 \frac{c(2\cdot4^m-1)}{2\cdot4^m-1}=\frac{4^{m+1}-1}{6\cdot4^m-3}>\frac23,
\]
the last inequality being \(3\cdot4^{m+1}-3>12\cdot4^m-6\). Thus
\(c(n)/n\) has no limit.

Let \(N\) be any null event. Since \(P(N\cup N_0)=0\), there is
\(\omega\notin N\cup N_0\). Put \(D_\omega=\{\omega_j:j\in J\}\), a
countable set. Because the coordinates of \(\omega\) are distinct,
\(\omega_j\in D_\omega\) exactly when \(j\in J\). So
\(A_{n,1}(\mathbf1_{D_\omega})(\omega)=c(n)/n\), which does not converge.
No null event has the property stated.

\emph{Item 5.} If \(\varphi_1,\varphi_2,\ldots\) satisfy the hypotheses
of \cref{thm:p15-exchangeable-averages}, each has a null event \(N_r\)
off which its averages converge to a fixed version of its limit. The
union \(\bigcup_rN_r\) is null, and off it every average converges.
\end{proof}
\end{annexitem}

\begin{annexitem}{annex:p15-w02c-p15-022-augmentation}{thm:p15-w02c-p15-022-augmentation}{Null augmentation enlarges the invariant, tail and kernel fields}
\begin{proof}
For any sub-sigma-field \(\mathcal G\), \(\mathcal N^*\subseteq\mathcal
G^{P}\): if \(P(A)=0\) take \(B=\varnothing\), and if \(P(A)=1\) take
\(B=\Omega\).

\emph{Item 1.} By \cref{thm:p15-hewitt-savage-zero-one}, every event of
\(\mathcal E\) has probability zero or one. If \(A\in\mathcal E^{P}\),
choose \(B\in\mathcal E\) with \(P(A\mathbin{\triangle}B)=0\); then
\(P(A)=P(B)\in\{0,1\}\), so \(A\in\mathcal N^*\). Together with the
first remark, \(\mathcal E^{P}=\mathcal N^*\). Since
\(\mathcal T\subseteq\mathcal E\),
\(\mathcal N^*\subseteq\mathcal T^{P}\subseteq\mathcal E^{P}=\mathcal N^*\).

\emph{Item 2.} By \cref{thm:p15-w02c-p15-022-field-strictness} item~2
with \(B=\{1\}\), \(A\in\mathcal E\setminus\mathcal T\). It is the union
over \(i\) of \(A^{(i)}=\{X_i=1,\ X_j=0\text{ for }j\neq i\}\). For
\(m\geq i\), \(A^{(i)}\subseteq\{X_i=1,\ X_j=0\text{ for }j\leq m,\
j\neq i\}\), which has probability \(2^{-m}\). So \(P(A^{(i)})=0\) and
\(P(A)=0\). Hence \(A\in\mathcal N^*=\mathcal T^{P}\) and
\(A\notin\mathcal T\).

\emph{Item 3.} Every event function of the kernel \(K\) is constant,
hence measurable with respect to \(\mathcal E\). For \(A\in\mathcal E\), \(P(A)\in\{0,1\}\) by
item~1. If \(P(A)=0\), \(E[\mathbf1_{\{X_1=1\}};A]=0=\tfrac12P(A)\); if
\(P(A)=1\), \(E[\mathbf1_{\{X_1=1\}};A]=P(X_1=1)=\tfrac12=\tfrac12P(A)\).
So \(K(\{1\})=1/2\) is a version of \(P(X_1=1\mid\mathcal E)\), and
\(K(\{0\})=1-K(\{1\})\), \(K(\varnothing)=0\), \(K(S)=1\) are versions of
the remaining event conditional probabilities. Hence \(K\) is a regular
conditional law.

For \(K'\), \(\bar Y\) takes values in \([0,1]\), so each \(K'_\omega\) is
a probability. As in \cref{annex:p15-w02c-p15-022-partial},
\(\bar Y\) is \(\mathcal T\)-measurable, hence \(\mathcal E\)-measurable;
so are \(K'(\{1\})=\bar Y\) and \(K'(\{0\})=1-\bar Y\). By
\cref{thm:p15-reverse-hewitt-savage} with \(k=1\) and the bounded
\(\varphi(x)=x\), \(S_n/n\to EX_1=1/2\) almost surely, so \(\bar Y=1/2\)
almost surely and each \(K'(G)\) equals \(K(G)\) almost surely. An
\(\mathcal E\)-measurable variable almost surely equal to a version of a
conditional probability is itself a version. So \(K'\) is also a regular
conditional law.

The event functions of \(K\) are constants, so
\(\sigma(K)=\{\varnothing,\Omega\}\). Next,
\(\sigma(K')=\sigma(\bar Y)\) contains \(\{\bar Y=1\}\), which contains
\((1,1,\ldots)\) and not \((0,0,\ldots)\); so
\(\sigma(K')\neq\sigma(K)\). Both fields are contained in \(\mathcal E\),
so their augmentations lie in \(\mathcal E^{P}=\mathcal N^*\); by the
first remark they contain \(\mathcal N^*\). Thus
\(\sigma(K)^{P}=\sigma(K')^{P}=\mathcal N^*=\mathcal E^{P}\). Finally the
event \(A\) of item~2 lies in \(\mathcal E\), contains
\((1,0,0,\ldots)\) and misses \((0,0,\ldots)\), so
\(A\notin\sigma(K)\) and \(\mathcal E\not\subseteq\sigma(K)\).
\end{proof}
\end{annexitem}

\begin{annexitem}{annex:p15-w02c-p15-022-cocountable}{thm:p15-w02c-p15-022-cocountable}{No countable determining class}
\begin{proof}
\emph{Item 1.} The definition is symmetric under complements. If every
member of a countable family is countable, the union is countable; if
some member has countable complement, the union contains it, and so has
countable complement. Every singleton is countable.

\emph{Item 2.} For each \(i\) let \(C_i=A_i\) if \(A_i\) is countable and
\(C_i=\mathbb R\setminus A_i\) otherwise. The union
\(U=\bigcup_iC_i\) is countable, so \(\mathbb R\setminus U\) is
uncountable; choose distinct \(x,y\) in it. If \(A_i\) is countable,
\(A_i\subseteq U\) contains neither point. Otherwise
\(\mathbb R\setminus A_i\subseteq U\), so \(A_i\) contains both.

(a) The class \(\mathcal M\) of sets \(H\subseteq\mathbb R\) with
\(x\in H\Leftrightarrow y\in H\) is closed under complements and countable
unions, so it is a sigma-field, and it contains every \(A_i\). Hence
\(\sigma(A_1,A_2,\ldots)\subseteq\mathcal M\). But
\(\{x\}\in\mathcal A\setminus\mathcal M\). So no countable family
generates \(\mathcal A\).

(b) The point masses \(\delta_x\) and \(\delta_y\) are countably additive
on every sigma-field, in particular on \(\mathcal A\). By the choice of
\(x,y\), \(\delta_x(A_i)=\delta_y(A_i)\) for every \(i\), while
\(\delta_x(\{x\})=1\neq0=\delta_y(\{x\})\). So no countable subclass of
\(\mathcal A\) determines probability measures on \(\mathcal A\). A
countable determining pi-system generating \(\mathcal A\), as required
in \cref{def:p15-directing-law}, is in particular such a subclass, so it
does not exist.

(c) Let \(\mathscr P\) be the probability measures on
\((\mathbb R,\mathcal A)\). The class of sets \(H\subseteq\mathscr P\)
with \(\delta_x\in H\Leftrightarrow\delta_y\in H\) is a sigma-field. It
contains every \(e_{A_i}^{-1}(G)\), \(G\) Borel, because
\(e_{A_i}(\delta_x)=e_{A_i}(\delta_y)\). It does not contain
\(e_{\{x\}}^{-1}(\{1\})\), which contains \(\delta_x\) and not
\(\delta_y\); that set belongs to the evaluation sigma-field. So the
evaluations \(e_{A_i}\) do not generate the evaluation sigma-field.

\emph{Item 3.} Suppose \(g\) were a Borel real coding of
\((\mathbb R,\mathcal A)\) onto a Borel set \(E\subseteq\mathbb R\), as in
\cref{def:p15-measurable-real-coding}. Measurability of \(g\) gives
\(g^{-1}(H)\in\mathcal A\) for every trace Borel set \(H\); measurability
of \(g^{-1}\) gives \(g(C)\) trace Borel for every \(C\in\mathcal A\), and
\(C=g^{-1}(g(C))\). So \(\mathcal A=g^{-1}(\mathcal B(E))\). The trace
field \(\mathcal B(E)\) is generated by the sets
\(E\cap(-\infty,q]\), \(q\) rational: the Borel sets \(H\) with
\(E\cap H\) in the field they generate form a sigma-field containing every
\((-\infty,q]\), hence every Borel set. For any class \(\mathcal C\) of
subsets of \(E\), \(g^{-1}(\sigma(\mathcal C))=\sigma(g^{-1}(\mathcal C))\):
the left side is a sigma-field containing \(g^{-1}(\mathcal C)\), and the
sets \(H\subseteq E\) with \(g^{-1}(H)\in\sigma(g^{-1}(\mathcal C))\)
form a sigma-field containing \(\mathcal C\). Hence \(\mathcal A\) would be
generated by the countable family \(g^{-1}(E\cap(-\infty,q])\),
contradicting item~2(a).
\end{proof}
\end{annexitem}

\begin{annexitem}{annex:p15-w02c-p15-022-urn-boundary}{thm:p15-w02c-p15-022-urn-boundary}{The finite-urn bound computed, vacuous and unattainable}
\begin{proof}
All laws in this proof live on a finite set of binary words, every
subset of which is an event. For a probability \(\lambda\) on \([0,1]\)
and a word \(w\in\{0,1\}^k\) with \(s\) ones,
\(M_\lambda^{(k)}(\{w\})=\int\theta^s(1-\theta)^{k-s}\,d\lambda(\theta)\);
the integrand is continuous and bounded, and the sum over all words is
\(\int1\,d\lambda=1\), so \(M_\lambda^{(k)}\) is a probability. Two facts
are used.

\emph{(T) Total variation on a finite set.} Let \(P,Q\) be
probabilities on a finite set \(W\) and \(A_+=\{w:P(w)>Q(w)\}\). For
every \(A\subseteq W\),
\(P(A)-Q(A)=\sum_{w\in A}(P(w)-Q(w))\leq\sum_{w\in A\cap A_+}(P(w)-Q(w))
\leq P(A_+)-Q(A_+)\), and, since both totals are \(1\),
\(Q(A)-P(A)=P(W\setminus A)-Q(W\setminus A)\leq P(A_+)-Q(A_+)\).
Equality holds at \(A=A_+\). Hence
\(\|P-Q\|_{\rm TV}=\sum_w(P(w)-Q(w))^+\).

\emph{(M) Images.} If \(f:W\to W'\), then for every \(B\subseteq W'\),
\(|P(f^{-1}B)-Q(f^{-1}B)|\leq\|P-Q\|_{\rm TV}\), so the image laws are
no farther apart than \(P\) and \(Q\).

Write \(z_1,\ldots,z_N\) for the population, with \(m\) ones, and recall
that \(H_{m,N}^{(k)}\) is the law of \((z_{J_1},\ldots,z_{J_k})\) for a
tuple \(J\) uniform over the \(N(N-1)\cdots(N-k+1)\) injective tuples in
\(\{1,\ldots,N\}\).

\emph{Item 1.} With \(k=2\) there are \(N(N-1)\) equally likely
injective pairs. Counting them,
\[
 H_{m,N}^{(2)}(11)=\frac{m(m-1)}{N(N-1)},\quad
 H_{m,N}^{(2)}(00)=\frac{(N-m)(N-m-1)}{N(N-1)},\quad
 H_{m,N}^{(2)}(01)=H_{m,N}^{(2)}(10)=\frac{m(N-m)}{N(N-1)}.
\]
Put \(c=m(N-m)/(N^2(N-1))=p(1-p)/(N-1)\). Then
\[
 \frac{m(m-1)}{N(N-1)}-\frac{m^2}{N^2}
 =\frac{m\bigl(N(m-1)-m(N-1)\bigr)}{N^2(N-1)}
 =\frac{m(m-N)}{N^2(N-1)}=-c,
\]
the same computation with \(N-m\) in place of \(m\) gives
\(H_{m,N}^{(2)}(00)-(1-p)^2=-c\), and
\[
 \frac{m(N-m)}{N(N-1)}-\frac{m(N-m)}{N^2}
 =\frac{m(N-m)\bigl(N-(N-1)\bigr)}{N^2(N-1)}=c
\]
for each of \(01\) and \(10\). By (T), \(d_{m,N}^{(2)}=2c\).
Since \(N^2-4m(N-m)=(N-2m)^2\geq0\), \(4m(N-m)\leq N^2\) with equality
exactly when \(N=2m\); dividing by \(2N^2(N-1)\) gives
\(2c\leq1/(2(N-1))\) with the same equality case. Next,
\(2c=1/N\) means \(2m(N-m)=N(N-1)\). As \(2m(N-m)\leq N^2/2\), this
forces \(N(N-1)\leq N^2/2\), that is \(N\leq2\); with \(N=2\) it reads
\(2m(2-m)=2\), so \(m=1\). Conversely \(N=2\), \(m=1\) gives
\(2c=1/2=1/N\). Finally
\(2c\big/(1/N)\leq N/(2(N-1))\), which is strictly less than \(1\) for
\(N\geq3\).

\emph{Item 2.} By item~1, \(d_{5,10}^{(2)}=2\cdot5\cdot5/(100\cdot9)=1/18\)
and \(d_{2,4}^{(2)}=2\cdot2\cdot2/(16\cdot3)=1/6\); the bounds are
\(2/20=1/10\) and \(2/8=1/4\). For \((m,N,k)=(2,4,3)\) there are \(24\)
injective triples. A word with one \(1\) is produced by \(2\) choices of
the one-entry at its \(1\) and \(2\cdot1\) ordered choices of the
zero-entries at its \(0\)s; a word with two \(1\)s by \(2\cdot1\) ordered
choices of the one-entries and \(2\) choices of the zero-entry. So each
of these six words has probability \(4/24=1/6\), and \(000\), \(111\)
have probability \(0\). The product law gives each of the eight words
probability \(1/8\). By (T),
\(d_{2,4}^{(3)}=6(1/6-1/8)=1/4\), against \(6/8=3/4\). The value
\(d_{2,4}^{(4)}=5/8\) is item~3 with \(\binom42 2^{-4}=3/8\), against
\(12/8=3/2\). For monotonicity, let \(f\) delete the last letter of a
word of length \(k+1\leq N\). The first \(k\) of the first \(k+1\) draws
are the first \(k\) draws, so \(f\) carries \(H_{m,N}^{(k+1)}\) to
\(H_{m,N}^{(k)}\), and it carries
\(\operatorname{Bernoulli}(p)^{\otimes(k+1)}\) to
\(\operatorname{Bernoulli}(p)^{\otimes k}\). By (M),
\(d_{m,N}^{(k)}\leq d_{m,N}^{(k+1)}\).

\emph{Item 3.} For \(k=N\), \(J\) is a uniform permutation. A word
\(w\) is produced exactly by the permutations that map the positions of
the \(1\)s of \(w\) onto the one-entries and the remaining positions
onto the zero-entries. If \(w\) has \(m\) ones there are \(m!(N-m)!\)
such permutations, and otherwise none. So
\(H_{m,N}^{(N)}(w)=1/\binom Nm\) on the \(\binom Nm\) words with \(m\)
ones and \(0\) elsewhere, while the product law gives each word with
\(m\) ones probability \(p^m(1-p)^{N-m}\). Put
\(b=\binom Nm p^m(1-p)^{N-m}\). By the binomial theorem \(b\) is one term
of \(\sum_{j=0}^N\binom Nj p^j(1-p)^{N-j}=1\), so \(b\leq1\) and
\(H_{m,N}^{(N)}(w)\geq p^m(1-p)^{N-m}\) on the words with \(m\) ones;
on every other word \(H_{m,N}^{(N)}\) is \(0\). By (T),
\(d_{m,N}^{(N)}=\sum_{w}\bigl(1/\binom Nm-p^m(1-p)^{N-m}\bigr)=1-b\), the
sum running over the words with \(m\) ones. If \(0<m<N\), then
\(0<p<1\), so \(b>0\), and the term \(j=0\neq m\) of the binomial sum
is \((1-p)^N>0\), so \(b<1\); hence \(0<d_{m,N}^{(N)}<1\). If \(m=0\),
every draw is \(0\), and \(\operatorname{Bernoulli}(0)^{\otimes k}\) is
the point mass at the all-zero word, which is also
\(H_{0,N}^{(k)}\); so \(d_{0,N}^{(k)}=0\) for every \(k\). The case
\(m=N\) is the same with ones. For \((m,N)=(5,10)\),
\(b=252\cdot2^{-10}=63/256\), so \(d_{5,10}^{(10)}=193/256\), while
\(10\cdot9/20=9/2\).

\emph{Item 4.} \(k(k-1)/(2N)\geq1\) exactly when \(k(k-1)\geq2N\), and
\(|P(A)-Q(A)|\leq1\) for any two probabilities and any event. For
\(N=4\) the values of \(k(k-1)\) for \(k=1,2,3,4\) are \(0,2,6,12\),
and only \(12\geq8\). For \(N=10\), \(k(k-1)\geq20\) holds at \(k=5\)
and, \(k(k-1)\) being increasing, at every larger \(k\), and fails at
\(k=4\), where \(k(k-1)=12\).

\emph{Item 5.} Let \(q_j=P(S_N=j)\). By
\cref{thm:p15-carnap-history-finite}, the law of \((X_1,\ldots,X_N)\) is
\(\sum_jq_jU_j\), where \(U_j\) is uniform over the words with \(j\)
ones. By the proof of item~3, \(U_j\) is the law of
\((z_{J_1},\ldots,z_{J_N})\) for a population with \(j\) ones and a
uniform permutation \(J\); each injective pair \((J_1,J_2)\) is the
initial pair of exactly \((N-2)!\) permutations, so \((J_1,J_2)\) is
uniform over injective pairs and the first two coordinates of \(U_j\)
have law \(H_{j,N}^{(2)}\). Hence
\(\mathcal L(X_1,X_2)=\sum_jq_jH_{j,N}^{(2)}\), and by definition
\(M_{\lambda_N}^{(2)}=\sum_jq_j\operatorname{Bernoulli}(j/N)^{\otimes2}\).
By the computation of item~1, for each word \(w\) of length \(2\),
\[
 \mathcal L(X_1,X_2)(w)-M_{\lambda_N}^{(2)}(w)
 =\pm\sum_jq_j\frac{(j/N)(1-j/N)}{N-1}=\pm\frac{E[V(1-V)]}{N-1},
\]
with the sign \(+\) for \(01\) and \(10\) and \(-\) for \(00\) and \(11\);
the sign does not depend on \(j\), so nothing cancels in the mixture.
By (T) the distance is \(2E[V(1-V)]/(N-1)\).

(a) For \(N=2\) and independent fair bits, \(V\) takes the values
\(0,1/2,1\) with probabilities \(1/4,1/2,1/4\), so
\(E[V(1-V)]=1/8\) and the distance is \(1/4\). The law of
\((X_1,X_2)\) is \(\operatorname{Bernoulli}(1/2)^{\otimes2}=M_{\delta_{1/2}}^{(2)}\).

(b) Here \(S_N=N/2\) almost surely, so \(\lambda_N=\delta_{1/2}\), and by
the preceding paragraph \(\mathcal L(X_1,X_2)=H_{N/2,N}^{(2)}\). Let
\(D=\{01,10\}\). By item~1,
\(H_{N/2,N}^{(2)}(D)=2(N/2)^2/(N(N-1))=N/(2(N-1))\). For any \(\lambda\),
\(M_\lambda^{(2)}(D)=2\int\theta(1-\theta)\,d\lambda(\theta)\leq1/2\),
because \(\theta(1-\theta)\leq1/4\) on \([0,1]\). Hence
\[
 \left\|\mathcal L(X_1,X_2)-M_\lambda^{(2)}\right\|_{\rm TV}
 \geq H_{N/2,N}^{(2)}(D)-M_\lambda^{(2)}(D)
 \geq\frac N{2(N-1)}-\frac12=\frac1{2(N-1)}.
\]
At \(\lambda=\delta_{1/2}\), \(M_\lambda^{(2)}=\operatorname{Bernoulli}(1/2)^{\otimes2}\),
and item~1 with \(m=N/2\) gives the distance \(1/(2(N-1))\). For \(N=2\)
this is \(1/2=2\cdot1/(2\cdot2)\), the bound at \(k=N=2\).
\end{proof}
\end{annexitem}
